\documentclass[a4paper,12pt]{article}
\usepackage[T2A]{fontenc}
\usepackage[utf8]{inputenc}
\usepackage[russian]{babel}
\usepackage{amsmath,amssymb,mathtools}
\usepackage{PTSans}
\renewcommand{\familydefault}{\sfdefault}
\usepackage{icomma}
\usepackage[a4paper,margin=19mm,headheight=25pt,headsep=9mm,footskip=14mm]{geometry}
\usepackage{microtype}
\usepackage{tikz}
\usetikzlibrary{calc,arrows.meta,positioning,shapes.geometric}
\usepackage{tabularx,booktabs,longtable}
\usepackage{enumitem}
\usepackage{xcolor}
\definecolor{accent}{HTML}{147D7E}
\definecolor{darkaccent}{HTML}{0B5253}
\definecolor{textgray}{HTML}{30383A}
\definecolor{softgray}{HTML}{6B7375}
\definecolor{rulegray}{HTML}{D7DEDE}
\usepackage[unicode=true,colorlinks=true,linkcolor=darkaccent,urlcolor=darkaccent]{hyperref}
\usepackage{parskip}
\usepackage{fancyhdr}
\usepackage{setspace}

\setstretch{1.10}
\setlength{\parindent}{0pt}
\setlength{\parskip}{5pt}
\AtBeginDocument{\setlength{\abovedisplayskip}{7pt}\setlength{\belowdisplayskip}{7pt}\setlength{\abovedisplayshortskip}{4pt}\setlength{\belowdisplayshortskip}{4pt}}
\setlist[itemize]{leftmargin=1.5em,itemsep=0.20em,topsep=0.25em}
\setlist[enumerate]{leftmargin=1.7em,itemsep=0.32em,topsep=0.28em}
\pagestyle{fancy}
\fancyhf{}
\fancyhead[L]{\small\color{softgray}Разложение уравнений на множители}
\fancyhead[R]{\small\color{softgray}08.10.26}
\fancyfoot[C]{\small\color{softgray}\thepage}
\renewcommand{\headrulewidth}{0.3pt}
\renewcommand{\headrule}{\hbox to\headwidth{\color{rulegray}\leaders\hrule height \headrulewidth\hfill}}
\newcommand{\sectionrule}{\par\vspace{-0.4em}\noindent{\color{rulegray}\rule{\linewidth}{0.6pt}}\par\vspace{0.25em}}
\newcommand{\key}[1]{\textcolor{darkaccent}{\textbf{#1}}}
\newcommand{\smallnote}[1]{{\small\color{softgray}#1}}
\tikzset{
  flowstart/.style={ellipse,draw=darkaccent,line width=0.85pt,minimum width=37mm,minimum height=10mm,align=center,fill=white,font=\small},
  flowaction/.style={rounded corners=2pt,draw=accent,line width=0.85pt,text width=63mm,minimum height=13mm,align=center,fill=white,font=\small,inner sep=3mm},
  flowdecision/.style={diamond,aspect=2.15,draw=accent,line width=0.85pt,text width=36mm,align=center,fill=white,font=\small,inner sep=1mm},
  flowarrow/.style={-{Stealth[length=2.6mm]},line width=0.85pt,draw=darkaccent}
}

\begin{document}
\color{textgray}
\begin{titlepage}
\thispagestyle{empty}
\centering
\vspace*{30mm}
{\large\color{softgray}Чек-лист}\par
\vspace{6mm}
{\Huge\bfseries\color{darkaccent}Решение уравнений\\[2mm]разложением на множители}\par
\vspace{5mm}
{\Large\color{accent}Группировка, общий множитель\\и разность квадратов}\par
\vspace{15mm}
\rule{0.68\linewidth}{0.7pt}\par
\vspace{9mm}
{\large Ярослав Верещагин}\par
\vspace{2mm}
{\normalsize 08.10.26}\par
\vfill
\begin{minipage}{0.83\linewidth}\centering\small
Цель: научиться узнавать структуру уравнения, получать произведение множителей,
находить все действительные корни и не терять решения при преобразованиях.
\end{minipage}
\vfill
{\small\color{softgray}Лёвин Артём Александрович эксклюзивно для Ярослава Верещагина}\par
\vspace{11mm}
\begin{tikzpicture}[remember picture,overlay]
\draw[accent,line width=1.15pt,opacity=0.5]
  ($(current page.south west)+(17mm,15mm)$)
  .. controls ($(current page.south west)+(62mm,31mm)$) and ($(current page.south east)+(-60mm,4mm)$) ..
  ($(current page.south east)+(-17mm,17mm)$);
\end{tikzpicture}
\end{titlepage}

\section*{1. Главный принцип: получить произведение, равное нулю}
\sectionrule
\begin{itemize}
  \item[$\square$] Переношу все слагаемые в одну часть: \(P(x)=0\).
  \item[$\square$] Ищу общий множитель, группировку, формулу сокращённого умножения.
  \item[$\square$] Получив \(A(x)B(x)=0\), решаю \(A(x)=0\) \textbf{или} \(B(x)=0\).
  \item[$\square$] Записываю \emph{все} действительные корни и проверяю ответ.
\end{itemize}

\smallnote{Количество слагаемых помогает выбрать приём: четыре слагаемых часто удобно группировать, но это не гарантия успешного разложения.}

\subsection*{Группировка: два последовательных вынесения}
Пример занятия:
\[
\begin{aligned}
x^3&=x^2-7x+7,\\
x^3-x^2+7x-7&=0,\\
(x^3-x^2)+(7x-7)&=0,\\
x^2(x-1)+7(x-1)&=0,\\
(x-1)(x^2+7)&=0.
\end{aligned}
\]
Так как \(x^2+7>0\) при любом действительном \(x\), остаётся \(x-1=0\). \key{Ответ: \(x=1\).}

\subsection*{Как правильно выносить отрицательный множитель}
\[
\begin{aligned}
x^3-3x^2-8x+24&=0,\\
x^2(x-3)-8(x-3)&=0,\\
(x-3)(x^2-8)&=0.
\end{aligned}
\]
Здесь второй общий множитель равен \(-8\), поэтому обе скобки становятся одинаковыми: \(x-3\).
\[
 x=3 \quad\text{или}\quad x^2=8\ \Longrightarrow\ x=\pm 2\sqrt2.
\]
\key{Ответ: \(x=3,\;x=-2\sqrt2,\;x=2\sqrt2\).}

\subsection*{Знаки при извлечении корня}
\[
\begin{array}{rcl@{\qquad}l}
x^2=9&\Longrightarrow&x=\pm3,&\text{два действительных корня};\\
x^2=0&\Longrightarrow&x=0,&\text{один корень};\\
x^2=-9&\Longrightarrow&\emptyset,&\text{нет действительных корней}.
\end{array}
\]

\clearpage
\section*{2. Общий множитель может быть целой скобкой}
\sectionrule
\subsection*{Сначала вынести \(x\), затем разложить трёхчлен}
\[
\begin{aligned}
x^3&=4x^2+5x,\\
x(x^2-4x-5)&=0,\\
x(x-5)(x+1)&=0.
\end{aligned}
\]
\key{Ответ: \(x=-1,\ 0,\ 5\).} Не делите уравнение на \(x\): такое деление потеряло бы корень \(x=0\).

\subsection*{Два общих множителя сразу}
\[
 (x-4)(x-5)(x-3)=(x-4)(x-5)(x-2).
\]
Переносим правую часть влево, выносим \emph{всё общее}:
\[
\begin{aligned}
(x-4)(x-5)\bigl[(x-3)-(x-2)\bigr]&=0,\\
-(x-4)(x-5)&=0.
\end{aligned}
\]
\key{Ответ: \(x=4\) или \(x=5\).} Скобки после знака минус обязательны:
\( -(x-2)=-x+2\).

\subsection*{Квадратный трёхчлен внутри более сложного уравнения}
Если \(ax^2+bx+c\) имеет корни \(x_1,x_2\) и \(a\ne0\), то
\[
 ax^2+bx+c=a(x-x_1)(x-x_2).
\]
Частный полезный случай: \(x^2+2x+1=(x+1)^2\). Например,
\[
\begin{aligned}
x(x^2+2x+1)&=2(x+1),\\
x(x+1)^2-2(x+1)&=0,\\
(x+1)\bigl[x(x+1)-2\bigr]&=0,\\
(x+1)(x+2)(x-1)&=0.
\end{aligned}
\]
\key{Ответ: \(x=-2,-1,1\).}

\subsection*{Ошибка, которую важно устранить}
При вынесении множителя из выражения \(AB-B\) остаётся \(B(A-1)\):
\[
 x(x-2)^2-(x-2)=(x-2)\bigl[x(x-2)-1\bigr].
\]
\key{Единица не исчезает.} После этого решают \(x-2=0\) и \(x^2-2x-1=0\), получая \(x=2\) и \(x=1\pm\sqrt2\).

\clearpage
\section*{3. Разность квадратов и четвёртых степеней}
\sectionrule
\subsection*{Квадраты скобок: часто быстрее, чем раскрытие}
Для любых действительных \(a,b\):
\[
 a^2-b^2=(a-b)(a+b),\qquad
 a^2=b^2\ \Longleftrightarrow\ a=b\ \text{или}\ a=-b.
\]
Пример занятия:
\[
\begin{aligned}
(2x-3)^2&=(1-2x)^2,\\
\bigl[(2x-3)-(1-2x)\bigr]\bigl[(2x-3)+(1-2x)\bigr]&=0,\\
(4x-4)(-2)&=0.
\end{aligned}
\]
\key{Ответ: \(x=1\).} Вторая скобка \(-2\) нулём не становится.

\subsection*{Четвёртая степень: дважды используем структуру квадрата}
\[
 a^4-b^4=(a^2-b^2)(a^2+b^2)=(a-b)(a+b)(a^2+b^2).
\]
Для действительных чисел \(a^2+b^2\ge0\). Оно равно нулю лишь при \(a=b=0\); этот случай уже учтён равенствами \(a=b\) и \(a=-b\). Поэтому
\[
 \boxed{a^4=b^4\ \Longleftrightarrow\ a=b\ \text{или}\ a=-b.}
\]
Например, в задаче \((2x-1)^4=(3-2x)^4\) разбираем два случая:
\[
\begin{aligned}
2x-1=3-2x&\ \Longrightarrow\ x=1;\\
2x-1=-(3-2x)&\ \Longrightarrow\ -1=-3\quad\text{(невозможно)}.
\end{aligned}
\]
\key{Ответ: \(x=1\).} Раскрывать скобки в четвёртую степень здесь не требуется.

\subsection*{Чек-лист типичных ошибок}
\begin{itemize}
\item[$\square$] Не забываю заменить \(=\) на \(=0\) перед применением свойства произведения.
\item[$\square$] В \((U-V)-(W-Z)\) раскрываю вторую скобку со сменой обоих знаков.
\item[$\square$] Если вынесение «съело всё второе слагаемое», оставляю \(1\) или \(-1\).
\item[$\square$] При \(x^2=c>0\) записываю \emph{два} знака перед корнем.
\item[$\square$] Не делю на выражение с \(x\), пока не проверил возможность его обнуления.
\item[$\square$] Если одна ветвь уравнения невозможна, записываю корни остальных ветвей.
\end{itemize}

\clearpage
\section*{4. Алгоритм: от уравнения к произведению}
\sectionrule
\vspace{5mm}
\begin{center}
\begin{tikzpicture}[node distance=8mm and 19mm]
\node[flowstart] (start) {Дано уравнение};
\node[flowaction,below=of start] (zero) {Перенести всё в одну часть:\\ \(P(x)=0\)};
\node[flowdecision,below=10mm of zero] (common) {Есть общий\\ множитель?};
\node[flowaction,below left=12mm and 5mm of common,text width=49mm] (yes) {Вынести \(x\), скобку\\ или несколько общих множителей};
\node[flowaction,below right=12mm and 5mm of common,text width=49mm] (no) {Попробовать группировку\\ или разность квадратов};
\node[flowaction,below=29mm of common,text width=72mm] (factor) {Разложить оставшиеся выражения\\ на множители, если требуется};
\node[flowaction,below=of factor,text width=72mm] (split) {Для каждого множителя записать\\ уравнение «множитель \(=0\)»};
\node[flowaction,below=of split,text width=72mm] (solve) {Решить полученные уравнения;\\ учесть \(x^2=c\) и невозможные случаи};
\node[flowstart,below=of solve] (end) {Записать все корни};
\draw[flowarrow] (start) -- (zero);
\draw[flowarrow] (zero) -- (common);
\draw[flowarrow] (common) -- node[near start,above left,font=\small]{да} (yes.north east);
\draw[flowarrow] (common) -- node[near start,above right,font=\small]{нет} (no.north west);
\draw[flowarrow] (yes.south) |- (factor.west);
\draw[flowarrow] (no.south) |- (factor.east);
\draw[flowarrow] (factor) -- (split);
\draw[flowarrow] (split) -- (solve);
\draw[flowarrow] (solve) -- (end);
\end{tikzpicture}
\end{center}

\clearpage
\section*{5. Тренировочный блок — Путь А}
\sectionrule
\textbf{10 заданий.} Решите уравнения над множеством действительных чисел. Там, где удобно, используйте группировку, вынесение множителя и формулу разности квадратов. Решения оформляйте последовательно; ответ записывайте в каждой задаче.

\begin{enumerate}[itemsep=1.15em,topsep=1em]
  \item \(x^3-x^2+9x-9=0.\)
  \item \(x^3-4x^2-9x+36=0.\)
  \item \(x^3=3x^2+4x.\)
  \item \((x-2)(x-5)(x+1)=(x-2)(x-5)(x-4).\)
  \item \(x(x^2-6x+9)=2(x-3).\)
  \item \(x(x^2+4x+4)=3(x+2).\)
  \item \((2x-5)^2=(3-2x)^2.\)
  \item \((3x+1)^2=(5-3x)^2.\)
  \item \((2x-3)^4=(x+4)^4.\)
  \item \((3x+2)^4=(5-3x)^4.\)
\end{enumerate}

\vspace{4mm}
\smallnote{В заданиях 5--6 сперва найдите квадрат двучлена. В заданиях 9--10 не раскрывайте четвёртые степени: используйте структуру разности квадратов.}

\clearpage
\section*{6. Ответы и контрольные разложения}
\sectionrule
\renewcommand{\arraystretch}{2.05}
\begin{longtable}{@{}p{10mm}p{0.42\linewidth}p{0.43\linewidth}@{}}
\toprule
\textbf{№} & \textbf{Удобное разложение / условие} & \textbf{Ответ}\\
\midrule
\endhead
1 & \((x-1)(x^2+9)=0\) & \(1\)\\
2 & \((x-4)(x-3)(x+3)=0\) & \(-3,\ 3,\ 4\)\\
3 & \(x(x-4)(x+1)=0\) & \(-1,\ 0,\ 4\)\\
4 & \(5(x-2)(x-5)=0\) & \(2,\ 5\)\\
5 & \((x-3)(x^2-3x-2)=0\) & \(3,\ \dfrac{3-\sqrt{17}}2,\ \dfrac{3+\sqrt{17}}2\)\\
6 & \((x+2)(x+3)(x-1)=0\) & \(-3,\ -2,\ 1\)\\
7 & \((4x-8)(-2)=0\) & \(2\)\\
8 & \((6x-4)\cdot6=0\) & \(\dfrac23\)\\
9 & \(2x-3=x+4\) или \(2x-3=-(x+4)\) & \(-\dfrac13,\ 7\)\\
10 & \(3x+2=5-3x\); вторая ветвь невозможна & \(\dfrac12\)\\
\bottomrule
\end{longtable}
\vspace{4mm}
\smallnote{Самопроверка: подстановка каждого найденного корня в исходное уравнение подтверждает равенство. Для уравнений вида $x^2=c$ при $c>0$ записывайте два корня.}
\end{document}